% RESEARCH_DRAFT_SHA256: df7b06a7d837ed84bcd52a7912f7c0668ffa628e97ccb890bfb3941811a4a637
% Layout and source-reference wrappers expanded for the web renderer.
\documentclass[11pt]{article}
\usepackage[margin=25mm]{geometry}
\usepackage{fontspec,unicode-math}
\setmainfont{Latin Modern Roman}
\setmathfont{Latin Modern Math}
\usepackage{amsmath,amssymb,mathtools,xurl,hyperref}
\hypersetup{colorlinks=true,urlcolor=blue,linkcolor=blue}
\setlength{\parindent}{0pt}
\setlength{\parskip}{5pt}
\setlength{\emergencystretch}{3em}
\newcommand{\sref}[2]{\hyperlink{source:#1:#2}{[S#2]}}
\newcommand{\cataloguescope}[1]{\paragraph{Source question.}#1}
\newcommand{\floor}[1]{\left\lfloor#1\right\rfloor}
\begin{document}

\section*{Bruhat path enumeration, cd reconstruction, and the explicit-output barrier}
\noindent catalogue entry; catalogue entry. Research attempt dated 11 October 2026.
\paragraph{Result in scope.}This manuscript contains a source-specific attempt and complete proofs of the stated auxiliary calculations. It makes no new solution or novelty claim for the entire target.
\section{Original problem and preserved definition}
The following definition and source list are reproduced from the repository. Any earlier inspection dates and reported literature claims in this preserved text describe that source record, not additional verification by this attempt.
% BEGIN PRESERVED DEFINITION
\subsection{Definitions and mathematical statement}

\paragraph{Shared notation.}Standard number systems are $\mathbb N,\mathbb Z,\mathbb Q,\mathbb R,\mathbb C$. Specialized notation and index ranges follow the cited source and the definitions below.


A Coxeter system is a group generated by involutions with pairwise product orders specified by its Coxeter matrix. $\ell(w)$ is the shortest generator-word length. The Bruhat interval $[u,v]$ consists of elements between $u$ and $v$ in Bruhat order.
\paragraph{Record scope and qualifications.}
The catalogue title is retained as metadata. Named certificates, contracts, bounds or frameworks in that title are not asserted as results of the original source extract.
The exact normalization of $P_{u,v}$, $a_i$, and cd-coefficients is source-specific; their asserted equivalence is retained as supplied.

Source-record limitation: The core interval invariant Puv and coefficient normalizations are not identified by the supplied extraction.
\paragraph{Statement recorded in the supplied source.}
Given a Coxeter system $(W,S)$ and elements $u\leq v$ with Coxeter lengths $\ell(u)$ and $\ell(v)$, what is the computational complexity of computing the Bruhat-interval quantity $P_{u,v}$, equivalently the coefficients $a_i(u,v)$, or all cd-coefficients $[w]_{u,v}$, as a function of the Coxeter matrix and the input lengths? \href{https://doi.org/10.4171/owr/2007/04}{[S2]}
\paragraph{Source status.}
The supplied catalogue labels this source record open; that metadata was checked against the supplied snapshot on 2026-10-11. The label is preserved for provenance and is not an independent verification of present open status or mathematical correctness. The catalogue literature assessment is dated 2026-08-21.
\subsection{Short English statement}
Compute the listed Bruhat-interval polynomial invariants efficiently as a function of the group and input lengths.
\subsection{Sources}
\begin{itemize}
\item[S1] Ulam AI and the UnsolvedMath Contributors, UnsolvedMath: A Curated Collection of Open Mathematics Problems, catalogue entry. Catalogue curation and metadata: CC BY 4.0. The source extract is retained; definitions, English summaries and qualifications are editorial additions.
\url{https://huggingface.co/datasets/ulamai/UnsolvedMath}.
\item[S2] Geometric and Topological Combinatorics (2007), Oberwolfach Reports, DOI 10.4171/owr/2007/04, pp. 195--272
\url{https://doi.org/10.4171/owr/2007/04}.
\item Source-specific Oberwolfach Reports licences, including CC BY-SA 4.0 where applicable, are retained; see THIRD\_PARTY\_NOTICES.md.
\end{itemize}
% END PRESERVED DEFINITION
\subsection{Verbatim supplied catalogue question}
Given a Coxeter system $(W,S)$ and elements $u\leq v$ with Coxeter lengths $\ell(u)$ and $\ell(v)$, what is the computational complexity of computing the Bruhat-interval quantity $P_{u,v}$, equivalently the coefficients $a_i(u,v)$, or all cd-coefficients $[w]_{u,v}$, as a function of the Coxeter matrix and the input lengths?

\section{Recovered source scope}
The primary report identifies $P_{u,v}$ as the Kazhdan--Lusztig polynomial.
It defines the complete $cd$-index by descent words of directed paths in the
reflection-labelled Bruhat graph. Problem 8 says computing $P$ is equivalent
to computing the $a_i$, and then asks about ``even'' all $cd$ coefficients;
it does not assert a two-way efficient equivalence between $P$ and the full
$cd$-index. This repairs context in our discussion without altering the
preserved catalogue wording. We do not rely on the report's displayed
coefficient conversion formulas, whose precise indexing would need further
checking, or claim to prove its complete-$cd$-index existence theorem.

\section{Attempt and results}
\paragraph{An exact algorithm once the labelled acyclic graph is supplied.}
Let a finite directed acyclic graph have source $u$, target $v\ne u$, and
ordered edge labels. For a path with labels $t_1,\ldots,t_k$, assign the
noncommutative word $w_1\cdots w_{k-1}$, where $w_i=a$ if
$t_i<t_{i+1}$ and $w_i=b$ otherwise, following a fixed declared descent
convention. Start at $u$ with one empty path. For each first edge, create
state (endpoint, last label, empty word) with count one. Process vertices in
a topological order. Extending a state across an outgoing edge of label $t$
appends $a$ or $b$ according to comparison with the stored last label, and
adds its count into the new state. Sum states at $v$ over their final labels.
Induction on path length proves correctness: each nonempty path has a unique
last edge and shorter prefix, and the extension records exactly its last
comparison. Thus every path contributes once and no extraneous path occurs.
It is an exact finite procedure for the supplied graph, with states potentially
exponential in the longest path length. It is not a polynomial algorithm for
constructing the graph from compressed Coxeter input.

\paragraph{A proved finite conversion from an $ab$ expansion.}
Put $c=a+b$ and $d=ab+ba$. Among $ab$ words of each fixed ordinary degree,
use lexicographic order with $a>b$. Expanding any $cd$ word, its leading
$ab$ word is obtained by replacing $c$ with $a$ and $d$ with $ab$, and its
coefficient is one. These leading words are distinct: parse a resulting word
from the right, pairing every $b$ with the immediately preceding $a$ as a
$d$, and reading any remaining $a$ as a $c$. This is a unique inverse parse
for words built from the tokens $a,ab$.
Thus the expansions of $cd$ words of a fixed weighted degree are linearly
independent. If an $ab$ polynomial is known to lie in their span, recover its
coefficients by leading-term subtraction: its largest remaining word must be
the largest leading word among the still-present $cd$ terms, since all other
terms of each expansion are smaller. Subtract its integer coefficient times
that expansion and continue. The process terminates in the finite word list.
This proves uniqueness and a finite reconstruction algorithm, conditional on
membership in the span. In the Bruhat application that membership is the
source's theorem, not a theorem re-established in this attempt.

\paragraph{Dense output alone can be exponential.}
Let $D$ be a requested maximum weighted degree, with $\deg c=1$, $\deg d=2$.
There are at least $2^{\lfloor D/2\rfloor}$ different $cd$ words of degree
$D$: independently choose the block $cc$ or $d$ in each of
$\lfloor D/2\rfloor$ positions and append $c$ if $D$ is odd.
A format explicitly listing the coefficient of every such word, even zeros,
therefore needs exponentially many entries. This does not bound the size of a
sparse representation, one selected coefficient, or the Kazhdan--Lusztig
polynomial. Consequently a polynomial-time statement must distinguish those
outputs rather than infer identical costs from a mathematical conversion.
For a simple path-count check, take the Boolean interval on three commuting
generators labelled $1<2<3$. Its directed paths are the six permutations.
Their comparison words are $aa$, $bb$, twice $ab$, and twice $ba$, so the
path polynomial is $a^2+2ab+2ba+b^2=c^2+d$ by direct expansion.
This is an illustrative small interval, not a general complexity estimate.

\section{Adversarial audit and remaining scope}
The audit checked path initialization (one edge contributes the empty word),
word order in the leading-term argument, output format, and the source's
weaker wording about all coefficients. Missing are efficient Coxeter word and
interval construction bounds in the declared encoding, coefficient bit-size
analysis, and a complexity classification for $P_{u,v}$ itself. No computation
beyond hand enumeration of the six paths was performed. The preserved source
limitation is clarified by actual source reading rather than silently replaced
with a different invariant.

\section{Process and contributions}
Mathematics and model-assisted adversarial self-audit: OpenAI gpt-6-astra, ultra reasoning, as dispatched for this research batch; exact serving revision not exposed. Codex (GPT-6; exact serving revision not exposed) materially assisted with orchestration and format review. No human mathematical authorship, mathematical edits, or independent human proof checking is asserted.
Contribution prompt A was applied to the full definition and supplied source row, then prompt B was applied adversarially to the derived results. The item-specific audit above records the scope checks and rejected extrapolations. Local assembly and structural checks do not establish mathematical correctness. The source catalogue status is preserved in TOP\_PROBLEM and is separate from this attempt status.
\paragraph{Integration review.} Codex (GPT-6; exact serving revision not exposed) reviewed the written arguments and expanded layout and source-reference wrappers. Repository metadata and the production web renderer were checked using Python 3.13 and JavaScript with MathJax 3.2.2. This model-assisted review is not independent human proof verification; no PDF compilation or formal proof checking is asserted.

\section{Completion Estimate}
\noindent\textbf{COMPLETION: 0\%}
The proved material is a rederivation of known or auxiliary facts rather than a substantiated new advance on the original target. The finite explicit-graph procedures and dense-output bound do not classify Kazhdan--Lusztig computation from the Coxeter matrix and element words or prove efficient equivalence of all output formats.
\section{References}
The complete original Sources subsection above is retained, including catalogue attribution and licence notices. The additional list records primary sources actually consulted for this attempt on 11 October 2026; uncited source links reproduced above were not all rechecked.
\begin{itemize}
\item Oberwolfach Report 4/2007, Billera, The complete cd-index of a Bruhat interval, printed pp.209--211; full relevant passage read from cached primary PDF text on 11 October 2026. \url{https://ems.press/content/serial-article-files/46090}.
\end{itemize}
\section{Final status}
\textbf{Unresolved in this attempt.} The finite explicit-graph procedures and dense-output bound do not classify Kazhdan--Lusztig computation from the Coxeter matrix and element words or prove efficient equivalence of all output formats.
\end{document}
